\documentclass[a4paper,12pt]{article}
\usepackage{iftex}
\ifPDFTeX
  \usepackage[T2A]{fontenc}
  \usepackage[utf8]{inputenc}
  \usepackage[russian]{babel}
\else
  \usepackage{fontspec}
  \setmainfont{DejaVu Sans}
  \usepackage{polyglossia}
  \setdefaultlanguage{russian}
\fi
\usepackage{amsmath,amssymb,amsthm,mathtools}
\usepackage{icomma}
\usepackage[a4paper,margin=1.65cm,headheight=16pt]{geometry}
\usepackage{microtype}
\usepackage{tikz}
\usetikzlibrary{positioning}
\usepackage{pgfplots}
\pgfplotsset{compat=1.18}
\usepackage{tabularx,booktabs,longtable}
\usepackage{enumitem}
\usepackage{hyperref}
\usepackage{parskip}
\usepackage{fancyhdr}
\usepackage{setspace}
\usepackage{xcolor}
\usepackage[most]{tcolorbox}
\usepackage{multicol}

\definecolor{accent}{HTML}{287F7A}
\definecolor{darkaccent}{HTML}{1C3151}
\definecolor{textgray}{HTML}{333333}
\definecolor{softteal}{HTML}{E4F3EF}
\definecolor{softcoral}{HTML}{FBE4DC}
\definecolor{coral}{HTML}{E77B5E}
\definecolor{softgold}{HTML}{F8EDCF}
\definecolor{softblue}{HTML}{E8EEF6}

\hypersetup{colorlinks=true,linkcolor=darkaccent,urlcolor=accent}
\setlist[itemize]{leftmargin=1.5em,itemsep=0.25em,topsep=0.35em}
\setlist[enumerate]{leftmargin=1.8em,itemsep=0.35em,topsep=0.35em}
\pagestyle{fancy}
\fancyhf{}
\lhead{Графики функций}
\rhead{Ксения Клыкова}
\cfoot{\thepage}
\renewcommand{\headrulewidth}{0.4pt}
\renewcommand{\today}{20 июля 2026 г.}

\newtcolorbox{keybox}[1][]{enhanced,breakable,colback=softteal,colframe=accent,boxrule=0.8pt,arc=3mm,left=4mm,right=4mm,top=3mm,bottom=3mm,title=#1,fonttitle=\bfseries,coltitle=darkaccent}
\newtcolorbox{warnbox}[1][]{enhanced,breakable,colback=softcoral,colframe=orange!70!red,boxrule=0.8pt,arc=3mm,left=4mm,right=4mm,top=3mm,bottom=3mm,title=#1,fonttitle=\bfseries,coltitle=darkaccent}
\newtcolorbox{examplebox}[1][]{enhanced,breakable,colback=softblue,colframe=darkaccent,boxrule=0.8pt,arc=3mm,left=4mm,right=4mm,top=3mm,bottom=3mm,title=#1,fonttitle=\bfseries,coltitle=darkaccent}
\newcommand{\point}[2]{\left(#1;#2\right)}

\begin{document}
\begin{titlepage}
\centering
\vspace*{2.3cm}
{\color{accent}\Large\bfseries ЧЕК-ЛИСТ\par}
\vspace{0.45cm}
{\color{darkaccent}\fontsize{31}{36}\selectfont\bfseries Графики функций:\\метод узловых точек\par}
\vspace{0.9cm}
{\large Прямая, парабола, гипербола и точки пересечения\par}
\vfill
{\Large\bfseries Лёвин Артём Александрович\\[0.2cm]эксклюзивно для Ксении Клыковой\par}
\vspace{0.8cm}
{\large \today\par}
\vfill
\begin{tikzpicture}[remember picture,overlay]
  \draw[accent,line width=2.2pt] (-7,-1.2) .. controls (-4,1.2) and (-1,-2.2) .. (1.2,0.2)
    .. controls (3.2,2.2) and (5.1,-0.8) .. (7,0.8);
  \fill[accent!35] (-6.1,-0.35) circle (3pt);
  \fill[accent!60] (0.2,-0.35) circle (3pt);
  \fill[accent!85] (6.2,0.25) circle (3pt);
\end{tikzpicture}
\end{titlepage}

\section*{1. Главная идея урока}
\begin{keybox}[Метод узловых точек]
Если формула графика содержит неизвестные коэффициенты, выберите на графике точки с точно читаемыми координатами и подставьте их координаты в формулу. Полученная система уравнений позволяет восстановить коэффициенты.
\end{keybox}

\begin{center}
\begin{tikzpicture}[node distance=8mm and 12mm,>=stealth,every node/.style={font=\small,align=center,text width=5.2cm,minimum height=1.05cm}]
\node[draw,rounded corners=3mm,fill=softgold] (a) {Неизвестные коэффициенты};
\node[draw,rounded corners=3mm,fill=softteal,right=of a] (b) {Столько же узловых точек};
\node[draw,rounded corners=3mm,fill=softblue,below=of a] (c) {Система уравнений};
\node[draw,rounded corners=3mm,fill=softcoral,right=of c] (d) {Формула и ответ};
\draw[->,thick,accent] (a)--(b); \draw[->,thick,accent] (b)--(d); \draw[->,thick,accent] (d)--(c);
\end{tikzpicture}
\end{center}

\subsection*{Сколько точек брать?}
\begin{itemize}
\item В формуле $y=kx+b$ неизвестны два коэффициента $k$ и $b$ --- нужны две точки.
\item В формуле $y=2x^2+bx+c$ неизвестны $b$ и $c$ --- также нужны две точки.
\item В формуле $y=\dfrac{k}{x}$ неизвестен только $k$ --- достаточно одной точки.
\item Выбирайте точки с целыми координатами; среди них удобнее те, чьи координаты ближе к нулю.
\end{itemize}

\section*{2. Как читать координаты и обозначения}
\begin{multicols}{2}
\begin{keybox}[Координаты точки]
Точка $A(x;y)$ записывается в порядке:
\[
\begin{aligned}
&\text{сначала абсцисса }x,\\
&\text{затем ордината }y.
\end{aligned}
\]
Для точки $A(-1;-3)$ подставляем $x=-1$ и $y=-3$.
\end{keybox}

\begin{keybox}[Два разных вопроса]
\[
f(-5)\quad\text{и}\quad f(x)=-5
\]
--- не одно и то же.
\begin{itemize}
\item $f(-5)$: подставить $x=-5$ и найти значение функции.
\item $f(x)=-5$: приравнять формулу к $-5$ и найти $x$.
\end{itemize}
\end{keybox}
\end{multicols}

\section*{3. Прямая: восстановление формулы}
\begin{examplebox}[Разобранный пример]
Пусть прямая $f(x)=kx+b$ проходит через точки $A(-1;-3)$ и $B(3;4)$. Найдём $f(-5)$.

Подставим координаты обеих точек:
\[
\begin{cases}
-3=-k+b,\\
4=3k+b.
\end{cases}
\]
Из первого уравнения $b=k-3$. Подставим во второе:
\[
3k+k-3=4 \quad\Longrightarrow\quad 4k=7 \quad\Longrightarrow\quad k=\frac74.
\]
Тогда
\[
b=\frac74-3=-\frac54.
\]
Формула прямой:
\[
f(x)=\frac74x-\frac54.
\]
Теперь вычислим значение функции:
\[
f(-5)=\frac74\cdot(-5)-\frac54=-\frac{40}{4}=-10.
\]
\textbf{Ответ: $-10$.}
\end{examplebox}

\begin{center}
\begin{tikzpicture}
\begin{axis}[
width=12.5cm,height=7.1cm,axis lines=middle,unit vector ratio=1 1,
xmin=-6,xmax=4.5,ymin=-12,ymax=6,samples=200,grid=both,
minor tick num=1,clip=false,xlabel={$x$},ylabel={$y$},
legend style={at={(0.03,0.97)},anchor=north west,draw=none,fill=white,fill opacity=.8,text opacity=1}]
\addplot[domain=-6:4.5,accent,very thick]{7*x/4-5/4};
\addlegendentry{$f(x)=\frac74x-\frac54$}
\addplot[only marks,mark=*,mark size=2.5pt,darkaccent] coordinates {(-1,-3) (3,4) (-5,-10)};
\node[anchor=south east] at (axis cs:-1,-3) {$A(-1;-3)$};
\node[anchor=south west] at (axis cs:3,4) {$B(3;4)$};
\node[anchor=north east] at (axis cs:-5,-10) {$(-5;-10)$};
\end{axis}
\end{tikzpicture}
\end{center}

\subsection*{Универсальный алгоритм для прямой}
\begin{enumerate}
\item Выберите две узловые точки.
\item Выпишите их координаты в порядке $(x;y)$.
\item Для каждой точки подставьте $x$ и $y$ в $y=kx+b$.
\item Решите систему относительно $k$ и $b$.
\item Запишите восстановленную формулу.
\item Выполните действие из вопроса: вычислите $f(a)$ или решите $f(x)=a$.
\item Проверьте подстановкой одной из исходных точек.
\end{enumerate}

\section*{4. Парабола: метод тот же}
\begin{examplebox}[Два неизвестных коэффициента]
Парабола имеет вид
\[
y=2x^2+bx+c
\]
и проходит через точки $A(-2;-2)$ и $B(1;1)$.

Подставляем координаты:
\[
\begin{cases}
-2=2\cdot(-2)^2-2b+c,\\
1=2\cdot1^2+b+c.
\end{cases}
\quad\Longrightarrow\quad
\begin{cases}
-2b+c=-10,\\
b+c=-1.
\end{cases}
\]
Вычитая второе уравнение из первого, получаем $-3b=-9$, поэтому $b=3$, а $c=-4$.

Искомая формула:
\[
y=2x^2+3x-4.
\]
\end{examplebox}

\begin{center}
\begin{tikzpicture}
\begin{axis}[
width=12cm,height=7.2cm,axis lines=middle,unit vector ratio=1 1,
xmin=-4,xmax=3,ymin=-7,ymax=9,samples=220,grid=both,minor tick num=1,clip=false,xlabel={$x$},ylabel={$y$}]
\addplot[domain=-4:3,coral!85!black,very thick]{2*x^2+3*x-4};
\addplot[only marks,mark=*,mark size=2.5pt,darkaccent] coordinates {(-2,-2) (1,1)};
\node[anchor=north east] at (axis cs:-2,-2) {$A(-2;-2)$};
\node[anchor=south west] at (axis cs:1,1) {$B(1;1)$};
\end{axis}
\end{tikzpicture}
\end{center}

\section*{5. Гипербола и пересечение графиков}
\subsection*{Как восстановить гиперболу}
Для функции
\[
f(x)=\frac{k}{x},\qquad x\ne0,
\]
достаточно одной точки. Например, если график проходит через $(1;2)$, то
\[
2=\frac{k}{1}\quad\Longrightarrow\quad k=2,
\]
следовательно, $f(x)=\dfrac2x$.

\subsection*{Условие пересечения}
В точке пересечения двух графиков их значения равны:
\[
f(x)=g(x).
\]
Это равенство действует именно для общих точек графиков.

\begin{examplebox}[Гипербола и прямая]
Пусть
\[
f(x)=\frac2x,\qquad g(x)=5x-9.
\]
Найдём абсциссы точек пересечения.
\[
\frac2x=5x-9,\qquad x\ne0.
\]
Домножим на $x$ с обязательным указанием ограничения:
\[
2=5x^2-9x \quad\Longrightarrow\quad 5x^2-9x-2=0.
\]
\[
D=(-9)^2-4\cdot5\cdot(-2)=121,
\]
\[
x_{1,2}=\frac{9\pm11}{10},\qquad x_1=2,\quad x_2=-\frac15=-0{,}2.
\]
Если точка $B$ расположена между $-1$ и $0$, выбираем $x_B=-0{,}2$.
\end{examplebox}

\begin{center}
\begin{tikzpicture}
\begin{axis}[
width=12.8cm,height=8cm,axis lines=middle,unit vector ratio=1 1,
xmin=-3,xmax=4,ymin=-12,ymax=10,samples=260,grid=both,minor tick num=1,clip=false,xlabel={$x$},ylabel={$y$},
legend style={at={(0.03,0.97)},anchor=north west,draw=none,fill=white,fill opacity=.85,text opacity=1}]
\addplot[domain=-3:-0.08,accent,very thick]{2/x};
\addplot[domain=0.08:4,accent,very thick]{2/x};
\addlegendentry{$f(x)=\frac2x$}
\addplot[domain=-3:4,coral!85!black,very thick]{5*x-9};
\addlegendentry{$g(x)=5x-9$}
\addplot[only marks,mark=*,mark size=2.7pt,darkaccent] coordinates {(2,1) (-0.2,-10)};
\node[anchor=south west] at (axis cs:2,1) {$A(2;1)$};
\node[anchor=north east] at (axis cs:-0.2,-10) {$B(-0{,}2;-10)$};
\end{axis}
\end{tikzpicture}
\end{center}

\section*{6. Типичные ошибки}
\begin{warnbox}[Контроль внимательности]
\begin{enumerate}
\item Перепутать порядок координат и вместо $(x;y)$ записать $(y;x)$.
\item Подставить только $x$, забыв, что $y=f(x)$ тоже берётся из точки.
\item Ошибиться при переносе: например, перенос $-4$ вправо даёт $+4$.
\item Неверно сложить дроби: $-\frac13+1=\frac23$, а не $-\frac43$.
\item Домножить уравнение с $\dfrac{k}{x}$ на $x$, не указав $x\ne0$.
\item В задаче о точке пересечения выписать оба корня, хотя рисунок однозначно указывает нужную точку.
\item Брать неудобные точки с дробными координатами, когда на графике есть целочисленные узлы.
\end{enumerate}
\end{warnbox}

\section*{7. Итоговый чек-лист}
\begin{keybox}[Перед записью ответа]
\begin{itemize}
\item Я определила вид функции и неизвестные коэффициенты.
\item Я взяла столько узловых точек, сколько коэффициентов неизвестно.
\item Я записала координаты в порядке $(x;y)$.
\item Я корректно подставила координаты и решила систему.
\item Я различила вопросы $f(a)$ и $f(x)=a$.
\item Для пересечения графиков я приравняла формулы.
\item При работе с гиперболой я записала $x\ne0$.
\item Я выбрала корень, соответствующий нужной точке на рисунке.
\item Я проверила знаки, дроби и одну исходную точку.
\end{itemize}
\end{keybox}

\clearpage
\section*{Тренировочный блок}
\textit{Решения здесь не приводятся. Сначала выполните задания самостоятельно, затем сверяйтесь с таблицей ответов на следующей странице.}

\subsection*{Средний уровень}
\begin{enumerate}
\item Прямая $f(x)=kx+b$ проходит через точки $(0;2)$ и $(3;8)$. Найдите $f(-2)$.
\item Прямая проходит через точки $(-2;5)$ и $(2;-3)$. Найдите $x$, если $f(x)=1$.
\item Прямая $f(x)=kx+b$ проходит через точки $(-1;-4)$ и $(2;5)$. Найдите $f(4)$.
\end{enumerate}

\subsection*{Уровень выше среднего}
\begin{enumerate}[resume]
\item Парабола $f(x)=x^2+bx+c$ проходит через точки $(0;-3)$ и $(2;3)$. Найдите $f(-2)$.
\item Парабола $f(x)=2x^2+bx+c$ проходит через точки $(-2;-2)$ и $(1;1)$. Найдите $f(0)$.
\item Гипербола $f(x)=\dfrac{k}{x}$ проходит через точку $(2;3)$. Найдите $f(-3)$.
\item Найдите обе абсциссы точек пересечения графиков $y=\dfrac2x$ и $y=5x-9$.
\item Найдите абсциссу точки пересечения графиков $y=\dfrac2x$ и $y=5x-9$, лежащую в интервале $(-1;0)$.
\item Гипербола $f(x)=\dfrac{k}{x}$ проходит через точку $(-4;2)$. Найдите $x$, если $f(x)=-1$.
\end{enumerate}

\subsection*{Сложный уровень}
\begin{enumerate}[resume]
\item Прямая $f(x)=kx+b$ проходит через точки $(1;4)$ и $(4;-2)$. Найдите абсциссу точки её пересечения с осью $Ox$ и объясните, почему в этой точке $f(x)=0$.
\end{enumerate}

\clearpage
\section*{Ответы для самопроверки}
\begin{center}
\renewcommand{\arraystretch}{1.35}
\begin{tabularx}{0.88\textwidth}{>{\bfseries}c X}
\toprule
№ & Ответ \\
\midrule
1 & $-2$ \\
2 & $0$ \\
3 & $11$ \\
4 & $-1$ \\
5 & $-4$ \\
6 & $-2$ \\
7 & $2$ и $-0{,}2$ \\
8 & $-0{,}2$ \\
9 & $8$ \\
10 & $3$; в точке пересечения с осью $Ox$ ордината равна нулю. \\
\bottomrule
\end{tabularx}
\end{center}

\vfill
\begin{keybox}[Как использовать ответы]
Если ответ не совпал, сначала проверьте порядок координат, знаки при переносе, арифметику дробей и соответствие выбранного корня положению точки на графике.
\end{keybox}

\end{document}
